%%%%%% ARMAN SHOKROLLAHI %%%%
\documentclass[a4paper]{arman-syllabus}
\usepackage{armansh}
\usepackage{armanshokrollahi}
\usepackage{geometry}
\geometry{inner=1.5cm,outer=1.5cm,top=2.5cm,bottom=2.5cm}
\newcommand{\courseterm}{پائیز ۱۳۹۳}
\newcommand{\coursetime}{چهارشنبه ۷ -- ۵ بعدازظهر}
\newcommand{\courseroom}{۲۱۰ ساختمان مهندسی}
\pagestyle{empty}
\usepackage{fancyhdr}
 \pagestyle{fancy} \fancyhf{}
% \fancyheadoffset[RO,LE]{30pt}
 \fancyhead[RO,LE]{ریاضیات گسسته 
 \hfill
  ۲ مهر ۱۳۹۳}
\fancyhead[LO]{\rightmark}
\fancyhead[RE]{\leftmark}
 \fancyfoot[C]{\thepage}
% \renewcommand\headrule
%  {{\color{blue}%
%    \hrule height 2pt
%           width\headwidth
%    \vspace{1pt}%
%    \hrule height 1pt
%           width\headwidth
%    \vspace{-4pt}}}

\usepackage{color,xecolor,graphicx}
\definecolor{darkblue}{rgb}{0.067,0.008,0.471}
\definecolor{darkgreen}{rgb}{0.008,0.871,0.067}
\definecolor{darkred}{rgb}{0.871,0.008, 0.067}
% \usepackage[colorlinks,citecolor=blue]{hyperref}
\usepackage{bbding, wasysym}
\usepackage{xepersian}

\setlatintextfont{Times New Roman}
\settextfont[Scale=1.2]{Amiri}

\begin{document}
\thispagestyle{empty}
\courseheader{ 
{به کلاس "ریاضیات گسسته`` خوش آمدید!}
%  \\

\section*{صفحه درس (\lr{AeLP})}  {\color{darkblue}{\url{http://yourwebsite.com/teaching}}}
\section*{ساعات پاسخگویی} 
\instructorhours{}
\section*{کتابهای درسی}  \coursetexts{}
نسخه‌ی الکترونیکی این کتابها برای دانلود و استفاده‌ی شما،  در سایت \lr{AeLP} موجود است.
\section*{اهداف درس} 
آیا کنجکاو هستید بدانید چه می‌دانید؟ چگونه باید بیان کرد مسئله‌ای حل می‌شود یا نمی‌شود؟ چه مسائل حل‌پذیری ارزش حل شدن ندارند؟ چه اتفاقاتی قرار است در آینده بیفتد؟ اگر این سؤالات از دید شما ارزش پرسیدن دارند، این کلاس همان جایی‌ست که باید دنبال جوابهای خود بگردید. \\
ریاضیات گسسته (یا ساختمانهای گسسته) مطالعه‌ی مباحثی از ریاضیات است که تغییر در آنها روان و هموار نیست و دارای مقادیر مختلف و ناهمگون می‌باشند، از قبیل اعداد صحیح، گرافها، یا عبارات منطقی. بدین ترتیب ریاضیات گسسته مباحث خود را از ریاضیات پیوسته (از قبیل حساب دیفرانسیل و آنالیز ریاضی) جدا می‌کند. ریاضیات گسسته با مجموعه‌های متناهی و نامتناهی اما شمارش‌پذیر، از اشیاء سروکار دارد. \\
 به طور کلی تعدادی از مباحثی که در مجموعه‌ی ریاضیات گسسته قرار میگیرند، عبارتند از: علوم کامپیوتر نظری، نظریه‌ی اطلاعات، منطق، نظریه‌ی مجموعه‌ها، ترکیبیات، نظریه‌ی گراف، احتمالات گسسته، نظریه‌ی اعداد، ساختارهای جبری از قبیل گروه‌ها، حلقه‌ها و مدولها،  نظریه‌ی جبری گراف، هندسه گسسته، هندسه ترکیبیاتی، تحقیق در عملیات، نظریه‌ی بازیها و مدل‌سازی داده‌ها. \\
 
در این درس ما زمان لازم برای یادگیری همه‌ی مباحث بالا را نداریم و فقط بخشی از آنها (که در بخش "سرفصل درس`` اشاره شده است) را مورد بررسی قرار میدهیم. 
در منطق ریاضیات، ما اثباتها را بیان می‌کنیم و یاد می‌گیریم که چگونه باید مطمئن شد آنچه می‌دانیم را واقعا می‌دانیم! اثبات چیست؟ آیا می‌توان یک نرم‌افزار را، بدون تست کردن آن، قابل اطمینان و معتبر دانست؟ \\
در بحث مجموعه‌ها و روابط نگاهی می‌اندازیم به گروهی از اشیاء که ویژگیهای مشابه دارند. زبانها چیستند؟ چگونه زمان اجرای یک الگوریتم را اندازه بگیریم؟ \\
در بحث گرافها، به مدلهای حمل‌و‌نقل و ارتباطات می‌پردازیم. چگونه بسته‌ها هدفهای خود را روی اینترنت می‌یابند؟ ساده‌ترین ماشین محاسباتی چه شکلی‌ست؟  "درخت``ها در واقع ساختارهایی هستند که برای سرعت طراحی شده‌اند. چه فاکتوری می‌تواند سریعترین محاسبات را سرعت ببخشد؟ چرا درختها را تقریباً همه جا در هوش مصنوعی و بیو‌انفورماتیک می‌توان یافت؟ \\
در بحث ترکیبیات و احتمالات، ما با امکان‌های موجود برای انجام کارها سروکار داریم. به چند حالت می‌توان افراد مختلف (زن و مرد) را در یک صف قرار داد؟ اگر بخواهیم زن‌ها و مردها یک‌در‌میان باشند چطور؟ در یک شبکه‌ی اجتماعی سه خانم و سه آقا وجود دارند، آیا می‌توان ادعا کرد که در این جمع یا سه نفر هستند که دوبه‌دو همدیگر را نمی‌شناسند یا سه نفر هستند که دو‌به‌دو همدیگر را می‌شناسند؟ \\
در بحث نظریه‌ی اعداد ما مقدماتی از محاسبات عددی را بیان می‌کنیم که کاربردهای حیاتی در فرستادن پیغام‌های رمزی دارند، بعنوان مثال شما پیغامی را با رمز برای دوستتان می‌فرستید، چطور میتوان مطمئن شد اگر وسط راه شخص سومی پیغام شما را دریافت کند قادر نخواهد بود آن را بخواند؟ قسمتی از این بحث را به متوقف کردن افراد بیگانه در فهمیدن پیغامهای خصوصی اختصاص خواهیم داد. این تکنیکها بر اساس درک درست ما از اعداد اول، فاکتورها، بخش‌پذیری اعداد و مدولها بنا شده‌اند. \\

در پایان این کلاس،‌ دانشجو باید قادر باشد:
\begin{itemize}
\item عملگرها را در سیستمهای مختلف عددی اجرا کند.
\item مسائل مختلف و عبارات محاوره‌ای را در قالب  نمادهای ریاضیاتی نمایش دهد و اثبات‌های معتبر را از نامعتبر تشخیص دهد.
\item اصول منطق ریاضی را یاد بگیرد و آنها را برای اثباتِ قدم‌به‌قدم نتایج ریاضیاتی و برنامه‌نویسی بکار گیرد.
\item مبانی منطق را در مدارهای دیجیتال به کار برد.
\item مبانی منطق گزاره‌ای را در استنتاج‌ها و بیان نتایج به درستی به کار برد.
\item انواع گرافها را بشناسد و طرق یافتن مسیرها و  مدارها را بلد باشد و درختهای مولد مینیمم تولید کند.
\item از مفاهیم "ترکیبیات`` و "جایگشت`` در یافتن احتمالات استفاده کند.
\item ساختارهای مهم جبری را بشناسد و ارتباط آنها با گراف‌ها و اعداد را بیان کند.
\item مقدماتی از الگوریتم‌ها مثل تعاریف و محاسبه‌ی پیچیدگی یک الگوریتم را بلد باشد.
\end{itemize}

اگر زمان اجازه دهد، می‌توانیم مفاهیم بیشتری را مورد مطالعه قرار دهیم، مفاهیمی از قبیل پیچیدگی الگوریتم‌ها، بهینه‌سازی خطی و ترکیبیاتی در مسائل برنامه‌ریزی و شبکه‌ها، بررسی گرافها به صورت الگوریتمی، و پوشاندن گرافها با استفاده از گرافهای کوچکتر.
\section*{سیستم نمره‌دهی} 
تمرینات و امتحانات کوتاه (٪۱۰), میان‌ترم اول (٪۲۵)، میان‌ترم دوم (٪۲۵)، پایان ترم (٪۴۰)

\section*{تاریخ‌های مهم}
\begin{tabular}{r|r}
  \hline
  % after \\: \hline or \cline{col1-col2} \cline{col3-col4} ...
   میان‌ترم 1 & ۲۱ آبان ۱۳۹۳ (۱۲ نوامبر، ۲۰۱۴) \\
  میان‌ترم 2 & ۲۶ آذر ۱۳۹۳ (۱۷ دسامبر, 2014)  \\
  پایان‌ترم & ۲۴ دی ۱۳۹۳  (۱۴ ژانویه ۲۰۱۵)  \\
  \hline
\end{tabular}  \\[.2cm]
تکرار امتحان برای دانشجویی که سر جلسه امتحان حاضر نشود، به هیچ وجه امکان‌پذیر نیست، مگر با هماهنگی قبلی و ارائه‌ی دلایل و مستندات قابل قبول.
\section*{قوانین درس}
\begin{itemize}
\item در سایت \lr{AeLP} ثبت نام کنید. بعد از ثبت نام،‌ بعد از تایید توسط استاد، شما به درس مربوطه لینک خواهید شد و آنگاه قادر خواهید بود محتویات درس را مشاهده نمایید.
\item هرهفته تعدادی تمرین به شما داده میشود که یک هفته فرصت دارید حل کنید (هفته‌ی بعد همان روز باید تحویل دهید). روزی که باید تمرینها را تحویل دهید، یک امتحان کوتاه ۱۰ دقیقه‌ای از شما گرفته میشود. امتحان بر مبنای سؤالات تمرین داده شده می‌باشد. شما میتوانید از تمرین و جوابهای خود سر جلسه امتحان کوتاه استفاده کنید ولی مجاز به استفاده از هیچگونه منبع دیگری نخواهید بود.
\item تمام جوابهای تمرینات باید تایپ‌ شده و در دو نسخه تحویل داده شوند: یک نسخه از آن سر کلاس همراه با امتحانِ کوناه و یک نسخه هم از طریق سایت \lr{AeLP} قبل از  اتمام زمانِ مورد نظر ذکر شده در برگه‌ی تمرینات. 
\item اگر دانشجو نسخه‌ی اولِ جواب تمرین‌های داده شده را قبل از زمان موردنظر، از طریق سایت ارسال نکند و نسخه‌ی دوم را همزمان با برگه‌ی امتحان کوتاه تحویل ندهد، امکان تحویل گرفتن آن در آینده نخواهد بود و نمره‌ی دانشجو در آن تمرین صفر می‌باشد.
\item شما برای حل تمرینات مجاز به مشورت با همکلاسی‌های خود  هستید، اما باید جوابها بر مبنای دانسته‌ها و یادگیری‌ها و به قلمِ خودِ شما باشد، و هیچگونه کپی‌برداری از جوابهای همکلاسی‌هایتان قابل قبول نخواهد بود. سایت \lr{AeLP} مجهز به سیستم تقلب‌ یاب  می‌باشد. شما مجاز به کپی‌برداری از اینترنت نیستید و آنچه به عنوان جواب تحویل می‌دهید صرفا باید نوشته‌ی خود شما بر اساس دانسته‌ها و یادگیری‌هایتان باشد.
\item از میان تمام تمرین‌های داده شده در طول ترم، کمترین نمره حذف و در نمره‌ی نهایی شما لحاظ نخواهد شد.
\item بهترین جواب برای هر تمرین، بنا به صلاحدید بنده، به عنوان جواب نهایی با اسم خود دانشجو در سایت گذاشته خواهد شد تا سایر دانشجوها استفاده کنند. در این صورت، دانشجوی موردنظر نمره‌ی تشویقی دریافت خواهد کرد.
\item صرفاً بدست آوردن جواب نهاییِ یک سؤال بدون ارائه توضیحات کامل و شرح راه‌حل، کاملا بی‌ارزش است. 
\item امتحانات میان‌ترم و پایان‌ترم به صورت کتاب-بسته (جزوه-بسته) می‌باشند و شما مجاز به استفاده از هیچگونه کتاب، جزوه و یا وسایل الکترونیکی از قبیل موبایل و تبلت نخواهید بود.
\end{itemize}


\section*{قوانین کلاس}
\begin{itemize}
\item  طبق قوانین دانشگاه، حضور مستمر در کلاسها اجباری می‌باشد و غیبت فقط با دلیل موجه قابل قبول است.
\item لطفاً موبایلهای خود را خاموش و یا در حالت بی‌صدا قرار دهید. اگر از لپ‌تاپ سرکلاس استفاده می‌کنید، صدای آن را قطع کرده و از کیبورد استفاده نکنید چون سروصدا باعث حواسپرتی سایر دانشجویان می‌شود.
\item اگر احیاناً سرکلاس غایب بودید، جزوه‌ی دوستانتان را کپی کنید. شما خود در قبال کلاس‌ها و مطالب گفته شده که غیبت کرده‌اید مسئولید. اگر جزوه‌ی دوستانتان را کپی کردید و مطالعه نمودید، آنگاه اگر سوالی داشتید میتوانید از بنده بپرسید.
\end{itemize}
\section*{دانشجویان ناتوان}
دانشجویان عزیزی که دارای ناتوانی فیزیکیِ خاصی هستند و نیاز به امکانات بیشتری دارند، لطفاً هرچه سریعتر (تا جلسه‌ی دوم کلاسها)‌ با من در میان بگذارند.
\section*{تقلب}
دانشجویانی که مرتکب تقلب و عدم صداقت آکادمیک (که توهین به استاد محسوب میشود) یا هرگونه تقلبی در راستای فریب استاد شوند، بسته‌به نوعِ‌  تقلب  در تمرینات یا امتحانات و بنابر تصمیم استاد، یکی از جریمه‌های زیر را پیش‌رویِ خود خواهند دید:
\begin{itemize}
\item انجام دادن مجدد آن تمرین 
\item  انجام دادن تمرین دیگری که ممکن است به آن دانشجو داده شود
\item گرفتن نمره‌ی صفر در آن تمرین یا امتحان
\item  گرفتن نمر‌ه‌ی صفر در کل درس به عنوان نمره‌ی پایان ترم
\end{itemize}

\section*{سرفصل درس}
مباحث زیر محتوای اصلی درس "ریاضیات گسسته`` در طول این ترم را تشکیل خواهند داد:
\begin{center} \begin{minipage}{5.8in}‎
‎\begin{flushleft}‎
مجموعه‌ها، ویژگیهای اعداد، توابع  \dotfill 
~فصل ۲ \lr{[R]}، ۱ جلسه \\
منطق ریاضی \dotfill  ~فصل ۱ \lr{[R]}، ۱ جلسه \\
 روشهای اثبات \dotfill  ~فصل ۱ \lr{[R]}، ۱ جلسه \\
   روشهای اثبات: کاربردها \dotfill ~فصول ۵،۳،۱   \lr{[R]}، ۲ جلسه \\
   جبر بولی، شبکه‌ها، و ساختارهای جبری \dotfill ~فصل ۱۲ \lr{[R]} و بخش ۴ \lr{[G]}، ۲ جلسه \\
   {\Large{\HandPencilLeft}} 
   میان‌ترم ۱ \dotfill  ~۱ ساعت (۲۱ آبان ۱۳۹۳) \\
مبانی نظریه اعداد و رمزگذاری \dotfill ~فصل ۴ \lr{[R]}، ۲ جلسه \\ 
  ترکیبیات و شمارش \dotfill  ~فصول ۹،۸،۷،۶ \lr{[R]} و فصول ۲،۱ \lr{[L]}، ۲ جلسه \\
  {\Large{\HandPencilLeft}} 
     میان‌ترم ۲ \dotfill   ~۱ ساعت (۲۶ آذر ۱۳۹۳) \\
نظریه‌ گراف و درختها  \dotfill  ~فصول ۱۱،۱۰ \lr{[R]} و فصول ۸،۷ \lr{[L]}، ۲ جلسه \\
 مدل‌بندی محاسبات، زبانها و گرامرها \dotfill ~فصل ۱۳ \lr{[R]}، ۱ جلسه
 \end{flushleft}‎
‎\end{minipage}‎
‎\end{center}‎


اگر زمان اضافه باقی ماند، سعی میکنیم مطالب زیر را هم به طور خیلی خلاصه بررسی کنیم:
\begin{center} 
\begin{minipage}{5.8in}
\begin{flushright}
نظریه کدگذاری و پیچیدگی محاسباتی \dotfill \\
آنالیز و پیچیدگی الگوریتم‌ها \dotfill \\
بهینه سازی خطی و ترکیبیاتی  \dotfill 
\end{flushright} \end{minipage} \end{center}

\section*{نکات}
\begin{itemize}
\item این درس دروازه‌ی ورود به مطالب مهم و بنیادی در علوم کامپیوتر و فنآوری اطلاعات است، بخصوص مباحث نظریه‌ گراف و نظریه اعداد. در یادگیری هرچه‌بیشتر آن کوشا باشید.
\item برای موفقیت در این درس باید دانشجو هفته‌ای حداقل  ۹--۶ ساعت زمان اختصاص دهد.
\item برای موفقیت در این درس و یادگیری مطالب آن باید صبر داشت و بارها و بارها مطالب را خواند.
‌\item امتحانات میان‌ترم و پایان‌ترم چیزی خارج از تمرینات هفتگی نیستند، پس این تمرینات را جدی بگیرید و کامل حل کنید. اگر مشکلی داشتید حضوراً بپرسید یا روی فروم‌های تعبیه شده در سایت \lr{AeLP} مطرح نمائید تا من یا دوستانتان در رفع مشکل کمک کنیم. سؤال شما ممکن است سؤال دانشجویان دیگر نیز باشد. هرگز از پرسیدن سؤالهایتان واهمه‌ای نداشته باشید چراکه هدف از آمدن به دانشگاه، یادگیری است.
\item نمره‌ی نهایی شما نشانگر میزان تلاش شما در این درس است، پس قابل بحث نخواهد بود.
\end{itemize}

\vspace{1cm}
\begin{center}
با آرزوی موفقیت و بهترین‌ها در این درس و در زندگی تحصیلی برای همه‌ی شما عزیزان {\Large \smiley}
\end{center}
%%%%% END
\end{document}